\documentclass[UTF8,a4paper,11pt]{ctexart}

\usepackage{amsmath,amssymb}
\usepackage{geometry}
\geometry{margin=2.5cm}
\usepackage{booktabs}
\usepackage{hyperref}
\usepackage{enumitem}

\begin{document}

在这张图右下角，你看到的是弹簧-质量系统的经典方程：

\[
m\ddot{x} = -kx
\]

其中：

\begin{itemize}
  \item $x$：位移（position），表示物体偏离平衡位置多少
  \item $\dot{x}$：速度（velocity），表示位移对时间的一阶导数
  \item $\ddot{x}$（读作 \emph{x double dot}）：加速度（acceleration），表示位移对时间的二阶导数
\end{itemize}

可以写成：

\[
\dot{x}=\frac{dx}{dt}
\]

\[
\ddot{x}=\frac{d^2x}{dt^2}
\]

或者用牛顿第二定律理解：

\[
F=ma
\]

所以

\[
m\ddot{x}
\]

其实就是质量 $m$ 乘以加速度 $a$。

\bigskip\noindent\rule{\linewidth}{0.4pt}\bigskip

这个弹簧系统里：

弹簧产生恢复力

\[
F=-kx
\]

（这叫胡克定律）

负号表示：

\begin{itemize}
  \item 往上拉（$x>0$）时，弹簧往下拉
  \item 往下压（$x<0$）时，弹簧往上推
\end{itemize}

因此：

\[
m\ddot{x}=-kx
\]

整理得

\[
\ddot{x}+\frac{k}{m}x=0
\]

这个方程的解是正弦波：

\[
x(t)=A\cos(\omega t+\phi)
\]

说明物体会一直上下振动。

\bigskip\noindent\rule{\linewidth}{0.4pt}\bigskip

结合你前面问的 \textbf{$G(s)$ 到底是什么}，这里其实就能看到控制理论的来源：

\begin{itemize}
  \item $x(t)$ 是时域里的位移
  \item $\ddot{x}(t)$ 是时域里的加速度
  \item 微分方程 $m\ddot{x}=-kx$ 描述系统运动规律
\end{itemize}

把这个方程做拉普拉斯变换后，就能得到系统传递函数 $G(s)$。

所以：

\begin{quote}
\textbf{$x$ 是系统的输出（位移），$\ddot{x}$ 是输出变化的速度再变化一次（加速度），而 $G(s)$ 是把输入力和输出位移联系起来的数学模型。}
\end{quote}

控制工程学的很多内容，本质上都是在研究这种弹簧、汽车悬架、电机、无人机等动态系统的运动规律。



\end{document}
